\documentclass[11pt]{article}
\usepackage[margin=25mm]{geometry}
\usepackage{fontspec,unicode-math}
\setmainfont{Latin Modern Roman}
\setmathfont{Latin Modern Math}
\usepackage{amsmath,amssymb,mathtools}
\usepackage{xurl,hyperref}
\hypersetup{colorlinks=true,urlcolor=blue,linkcolor=blue}
\setcounter{secnumdepth}{0}
\setlength{\parindent}{0pt}
\setlength{\parskip}{5pt}
\setlength{\emergencystretch}{3em}
\newcommand{\sref}[2]{\hyperlink{source:#1:#2}{[S#2]}}
\newcommand{\cataloguescope}[1]{\paragraph{Source question.}#1}
\newcommand{\floor}[1]{\left\lfloor#1\right\rfloor}
\begin{document}
\section{Orbit-Closure Obstructions for Permanent–Determinant Separation}\label{statement}

\subsection{Definitions and mathematical statement}

\paragraph{Shared notation.}Standard number systems are denoted by $\mathbb N,\mathbb Z,\mathbb Q,\mathbb R,\mathbb C$; index ranges, computational models and specialized notation follow the cited problem source and the definitions below.


The determinant and padded permanent are homogeneous polynomials acted on by linear variable changes. Equivariance of a ring map means it commutes with this group action; a surjection is a map onto the whole target ring.

For a group action on a vector space, $Gp$ is the orbit of $p$ and $\overline{Gp}$ is its algebraic closure. Its coordinate ring consists of polynomial functions modulo those vanishing on the orbit closure. Representation multiplicities count copies of a given irreducible representation in a graded component.
\paragraph{Statement recorded in the supplied source.}
For every positive polynomial $p$ and every $m_0\in\mathbb{N}$, does there exist $m\geq m_0$ such that no equivariant surjection exists from the coordinate ring of the determinant orbit closure to that of the padded permanent orbit closure? \sref{statement}{2}
\paragraph{Scope qualification.}
The padding, dimension polynomial $p$, and group representation are the conventions of the source.
\paragraph{Source status.}
The supplied catalogue labels this record partially solved. Its source-status metadata was checked against the supplied snapshot on 2026-10-11. This import does not independently establish current open status or mathematical correctness. The catalogue literature assessment is dated 2026-08-21.
\subsection{Short English statement}
Rule out the stated equivariant ring surjections between determinant and padded-permanent orbit closures at arbitrarily large sizes.
\subsection{Sources}
\begin{itemize}
\item[\hypertarget{source:statement:1}{S1}] Ulam AI and the UnsolvedMath Contributors, UnsolvedMath: A Curated Collection of Open Mathematics Problems, catalogue entry. Catalogue curation and metadata: CC BY 4.0. The supplied source question is retained; definitions and the English summary are editorial additions.
\url{https://huggingface.co/datasets/ulamai/UnsolvedMath}.
\item[\hypertarget{source:statement:2}{S2}] Complexity Theory (2010), Oberwolfach Reports, DOI 10.4171/owr/2009/52, pp. 2787--2850
\url{https://doi.org/10.4171/owr/2009/52}.
\item[\hypertarget{source:statement:3}{S3}] permanent versus determinant survey
\url{https://arxiv.org/abs/1511.08113}.
\textit{Underlying Oberwolfach Reports material retains its source-specific licence, including CC BY-SA 4.0 where applicable; see THIRD\_PARTY\_NOTICES.md.}
\end{itemize}
\paragraph{Preparation and review.}
Original question attributed to the cited source and UnsolvedMath contributors. Codex (GPT-6; exact serving revision not exposed) assisted with definitions, English summaries, source inspection and formatting. No new human mathematical authorship or independent proof verification is asserted. Catalogue identity checked against the complete supplied registry and existing repository statements; known flagged repetitions and obvious identical conjectures excluded. Full mathematical-equivalence/literature review has not been performed. Statement preparation and final review followed contribution prompts C and D; automated checks verify structure and identity, not open status.
\label{end:statement}
\end{document}
